
\documentclass{article}
\usepackage{colortbl}
\usepackage{makecell}
\usepackage{multirow}
\usepackage{supertabular}

\begin{document}

\newcounter{utterance}

\centering \large Interaction Transcript for game `referencegame', experiment `number\_grids', episode 14 with reflexion{-}Qwen3.5{-}9B.
\vspace{24pt}

{ \footnotesize  \setcounter{utterance}{1}
\setlength{\tabcolsep}{0pt}
\begin{supertabular}{c@{$\;$}|p{.15\linewidth}@{}p{.15\linewidth}p{.15\linewidth}p{.15\linewidth}p{.15\linewidth}p{.15\linewidth}}
    \# & $\;$A & \multicolumn{4}{c}{Game Master} & $\;\:$B\\
    \hline

    \theutterance \stepcounter{utterance}  
    & & \multicolumn{4}{p{0.6\linewidth}}{
        \cellcolor[rgb]{0.9,0.9,0.9}{
            \makecell[{{p{\linewidth}}}]{
                \texttt{\tiny{[P1$\langle$GM]}}
                \texttt{You are given three grids, where each of them is 5 by 5 in size.} \\
\texttt{Grids have empty cells marked with "▢" and filled cells marked with "X".} \\
\texttt{Your task is to generate a referring expression that best describes the target grid while distinguishing it from the two other distractor grids.} \\
\texttt{The first grid is the target grid, and the following two grids are the distractors.} \\
\\ 
\texttt{Target grid:} \\
\\ 
\texttt{□ X X X X} \\
\texttt{□ X □ □ X} \\
\texttt{□ X □ □ X} \\
\texttt{□ X □ □ X} \\
\texttt{□ X X X X} \\
\\ 
\texttt{Distractor grid 1:} \\
\\ 
\texttt{□ X X X X} \\
\texttt{□ X □ □ X} \\
\texttt{□ X X X X} \\
\texttt{□ X □ □ X} \\
\texttt{□ X X X X} \\
\\ 
\texttt{Distractor grid 2:} \\
\\ 
\texttt{□ X X X X} \\
\texttt{□ X □ □ X} \\
\texttt{□ X X X X} \\
\texttt{□ □ □ □ □} \\
\texttt{□ X X X X} \\
\\ 
\texttt{Instruction: Describe the target grid.} \\
\texttt{Generate the referring expression starting with the tag "Expression: " for the given target grid. Omit any other text.} \\
            }
        }
    }
    & \\ \\

    \theutterance \stepcounter{utterance}  
    & \multicolumn{4}{p{0.6\linewidth}}{
        \cellcolor[rgb]{0.9,0.9,0.9}{
            \makecell[{{p{\linewidth}}}]{
                \texttt{\tiny{[P1$\rangle$GM]}}
                \texttt{Expression: The grid with a vertical line of filled cells in the second column and a horizontal line of filled cells in the first and last rows, forming a complete rectangular border of filled cells on the left, right, top, and bottom edges, with the interior cells being empty except for the second column.} \\
            }
        }
    }
    & & \\ \\

    \theutterance \stepcounter{utterance}  
    & & & \multicolumn{2}{p{0.3\linewidth}}{
        \cellcolor[rgb]{0.9,0.9,0.9}{
            \makecell[{{p{\linewidth}}}]{
                \texttt{\tiny{[GM$|$GM]}}
                \texttt{Expression: The grid with a vertical line of filled cells in the second column and a horizontal line of filled cells in the first and last rows, forming a complete rectangular border of filled cells on the left, right, top, and bottom edges, with the interior cells being empty except for the second column.} \\
            }
        }
    }
    & & \\ \\

    \theutterance \stepcounter{utterance}  
    & & \multicolumn{4}{p{0.6\linewidth}}{
        \cellcolor[rgb]{0.9,0.9,0.9}{
            \makecell[{{p{\linewidth}}}]{
                \texttt{\tiny{[GM$\rangle$P2]}}
                \texttt{You are given three grids, where each of them is 5 by 5 in size.} \\
\texttt{Grids have empty cells marked with "▢" and filled cells marked with "X".} \\
\texttt{You are also given a referring expression that describes one of the given grids.} \\
\texttt{Your task is to select the grid that matches the given referring expression.} \\
\texttt{Generate only the number (in text) of the grid that matches the given expression by selecting first, second, or third.} \\
\\ 
\texttt{First:} \\
\\ 
\texttt{□ X X X X} \\
\texttt{□ X □ □ X} \\
\texttt{□ X X X X} \\
\texttt{□ □ □ □ □} \\
\texttt{□ X X X X} \\
\\ 
\texttt{Second:} \\
\\ 
\texttt{□ X X X X} \\
\texttt{□ X □ □ X} \\
\texttt{□ X X X X} \\
\texttt{□ X □ □ X} \\
\texttt{□ X X X X} \\
\\ 
\texttt{Third:} \\
\\ 
\texttt{□ X X X X} \\
\texttt{□ X □ □ X} \\
\texttt{□ X □ □ X} \\
\texttt{□ X □ □ X} \\
\texttt{□ X X X X} \\
\\ 
\texttt{Expression: The grid with a vertical line of filled cells in the second column and a horizontal line of filled cells in the first and last rows, forming a complete rectangular border of filled cells on the left, right, top, and bottom edges, with the interior cells being empty except for the second column.} \\
\texttt{Question: Which grid does the expression refer to?} \\
\texttt{Start with the tag "Answer: ", followed by your selection. Omit any other text.} \\
            }
        }
    }
    & \\ \\

    \theutterance \stepcounter{utterance}  
    & & & \multicolumn{4}{p{0.6\linewidth}}{
        \cellcolor[rgb]{0.9,0.9,0.9}{
            \makecell[{{p{\linewidth}}}]{
                \texttt{\tiny{[GM$\langle$P2]}}
                \texttt{Answer: Second} \\
            }
        }
    }
     \\ \\

    \theutterance \stepcounter{utterance}  
    & & & \multicolumn{2}{p{0.3\linewidth}}{
        \cellcolor[rgb]{0.9,0.9,0.9}{
            \makecell[{{p{\linewidth}}}]{
                \texttt{\tiny{[GM$|$GM]}}
                \texttt{second} \\
            }
        }
    }
    & & \\ \\

    \theutterance \stepcounter{utterance}  
    & \multicolumn{4}{p{0.6\linewidth}}{
        \cellcolor[rgb]{0.95,0.95,0.95}{
            \makecell[{{p{\linewidth}}}]{
                \texttt{\tiny{[P2$\rightarrow$P1]}}
                \texttt{Post{-}episode Reflexion (not sent to GM):} \\
\texttt{New plan: My previous description was too vague and failed to explicitly distinguish the target from the distractors by pinpointing the exact row differences. Specifically, I mentioned a "complete rectangular border" which is true for the target, but I did not clearly state that the distractors have a "broken" bottom border (missing the bottom{-}left cell) or a "broken" middle section in the third row. Next time, I will construct the expression by first confirming the shared features (the vertical line in column 2 and the top/bottom rows in the target) and then immediately contrasting them with the specific missing cells in the distractors: explicitly stating that the target has filled cells at (row 3, col 1) and (row 4, col 1) which are empty in both distractors, and that the target has a solid bottom row while Distractor 1 has a gap at the bottom{-}left and Distractor 2 has a gap at the bottom{-}left and a different middle pattern. I will ensure the phrase "unlike the distractors" is followed by the precise coordinates or relative positions of the differing cells.} \\
            }
        }
    }
    & & \\ \\

    \theutterance \stepcounter{utterance}  
    & & & \multicolumn{4}{p{0.6\linewidth}}{
        \cellcolor[rgb]{0.95,0.95,0.95}{
            \makecell[{{p{\linewidth}}}]{
                \texttt{\tiny{[P1$\rightarrow$P2]}}
                \texttt{Post{-}episode Reflexion (not sent to GM):} \\
\texttt{New plan: I incorrectly identified the "Second" grid because I focused on the vertical line in the second column but failed to verify the horizontal border requirement. The expression demands a complete rectangular border (filled first and last rows) with an empty interior except for the second column. The "Second" grid has empty cells in the first and last rows (specifically columns 3 and 4), violating the border condition. The "Third" grid, however, has filled cells in the entire first and last rows, creating the required complete border, while maintaining the vertical line in the second column and an empty interior elsewhere. Next time, I will strictly check for the presence of filled cells in all four corners and the entirety of the first and last rows before confirming a match.} \\
\texttt{Answer: Third} \\
            }
        }
    }
     \\ \\

\end{supertabular}
}

\end{document}
